% ============================================================
%  ch14 —— §10、§11（原文 p.52–p.55）
% ============================================================
\routenote{p.52--p.55}

\section{附录：拟极小的组合论}

本节假设 $G$ 几乎是单的且秩至少为 $2$。

设 $\theta$ 是最高根（它是长根），$\theta^\vee$ 是对应的余根（它是短余根）。根集 $\Phi$ 被划分为
\[
\Phi^\theta_n=\{\alpha\in\Phi\mid\langle\alpha,\theta^\vee\rangle=n\},\qquad n=0,\pm1,\pm2,
\]
其中 $\Phi^\theta_{\pm2}=\{\pm\theta\}$。$\Phi^\theta_1$ 中的根成对配为 $(\alpha,\beta)$，满足 $\alpha+\beta=\theta$。

设 $V_{\theta^\vee}$ 是 $\ckG$ 的以 $\theta^\vee$ 为最高权的不可约表示。$V_{\theta^\vee}$ 的非零权正是 $\ckG$ 的短根，每个重数为 $1$。设 $\{e=\sum_ix_i,\ 2\rho,\ f=\sum_iy_i\}\in\ckg$ 是主 $\mathfrak{sl}_2$-三元组，其中 $x_i\in\ckg_{\alpha_i^\vee}$、$y_i\in\ckg_{-\alpha_i^\vee}$ 非零，$2\rho\in\ck t$ 是 $\ckG$ 的正余根之和。

\begin{lemma}
下列数相等：
\begin{enumerate}[label=(\arabic*), leftmargin=1.9em, itemsep=.2em]
\item $\dim V_{\theta^\vee}(0)$，其中 $V_{\theta^\vee}(0)$ 是 $2\rho$-作用下的零权空间，也是 $\check T$-作用下的零权空间；
\item $\dim V^e_{\theta^\vee}$，其中 $V^e_{\theta^\vee}=\ker(\operatorname{ad}(e)|_{V_{\theta^\vee}})$；
\item $\ckG$ 的短单根数（短秩 $r_s(\ckG)$）；
\item $\#(W\theta^\vee)/h=\#\{\ckG\ \text{的短根}\}/h$（这里 $h$ 是 $W$ 的 Coxeter 数）。
\end{enumerate}
\end{lemma}

\begin{proof}
在主 $\mathfrak{sl}_2$-作用下，$V_{\theta^\vee}$ 可分解为 $\mathfrak{sl}_2$ 的不可约表示之和：
\[
V_{\theta^\vee}=\bigoplus_{i=1}^{r_s(\ckG)}\operatorname{Sym}^{2\ell_i}(\mathrm{Std}),
\]
其中 $\mathrm{Std}$ 是 $\mathfrak{sl}_2$ 的二维表示。由于 $2\rho$-作用在 $V_{\theta^\vee}$ 上的权都是偶数，$V_{\theta^\vee}$ 中只出现 $\mathrm{Std}$ 的偶对称幂。每个 $\operatorname{Sym}^{2\ell_i}(\mathrm{Std})$ 对 $V_{\theta^\vee}(0)$ 与 $V^e_{\theta^\vee}$ 各贡献 $1$ 维，故
\[
r_s(\ckG)=\dim V_{\theta^\vee}(0)=\dim V^e_{\theta^\vee}.
\]
这证明了 (1) 与 (2) 中的数相等。

设 $V_{\theta^\vee}(n)$ 是 $\rho$-作用的权 $n$-特征空间，则
\[
V_{\theta^\vee}(n)=\bigoplus_{\alpha^\vee\ \text{短},\ \langle\rho,\alpha^\vee\rangle=n}V_{\theta^\vee}(\alpha^\vee).
\]
特别地，$\dim V_{\theta^\vee}(1)$ 是 $\ckG$ 的短单根数。映射 $e:V_{\theta^\vee}(0)\to V_{\theta^\vee}(1)$ 显然满。它也是单的：若 $v\in V_{\theta^\vee}(0)$ 且 $ev=0$，则对一切单根 $\alpha_i^\vee$ 有 $\ckg_{\alpha_i^\vee}v=0$，即 $v$ 是最高权向量，矛盾。这证明了
\begin{equation}
\dim V_{\theta^\vee}(0)=\dim V_{\theta^\vee}(1)=r_s(\ckG).
\end{equation}

还需证明 (4) 与其余各数相同。论证与 [28, §6.7] 类似，那里 Kostant 考虑的是伴随表示而非 $V_{\theta^\vee}$。这里只给梗概。设 $\zeta$ 是本原 $h$ 次单位根，$P=\rho(\zeta)\in\ckG$。设 $\gamma^\vee$ 是 $\ckG$ 的最高根，$z=e+x_{-\gamma^\vee}$。则 $z$ 是 $\ckg$ 中的正则半单元，且 $\operatorname{Ad}(P)z=\zeta z$。设 $\ckG_z$ 是 $z$ 的中心化子（它是极大环面），则 $P\in N_{\ckG}(\ckG_z)$，其在 Weyl 群中的像是 Coxeter 元。取 $\operatorname{Sym}^{2\ell_i}(\mathrm{Std})\subset V_{\theta^\vee}$ 中最高权线的一组基 $u_i$，则存在唯一的 $v_i\in V_{\theta^\vee}(\ell_i-h)$ 使 $u_i+v_i\in V^z_{\theta^\vee}$（即 $z$-作用在 $V_{\theta^\vee}$ 上的核）。于是 $\{u_i+v_i\}_{1\le i\le r_s(\ckG)}$ 构成 $V^z_{\theta^\vee}$ 的基，在 $\operatorname{Ad}(P)$ 作用下的特征值为 $\zeta^{\ell_i}$。

$\ckG$ 的表示 $V_{\theta^\vee}$ 显然自对偶。按 $\ckg^z$ 的零权与非零权，可把 $V_{\theta^\vee}$ 写成 $N_{\ckG}(\ckG_z)$-模的直和 $V_{\theta^\vee}=V^z_{\theta^\vee}\oplus V'$。$\ckG$-模的任何自对偶 $V_{\theta^\vee}\cong V^*_{\theta^\vee}$ 必然限制为 $N_{\ckG}(\ckG_z)$-模的自对偶 $V^z_{\theta^\vee}\cong(V^z_{\theta^\vee})^*$。因此 $\operatorname{Ad}(P)$ 在 $V^z_{\theta^\vee}$ 上的特征值在求逆下不变：即多重集 $\{\zeta^{\ell_i}\}_{1\le i\le r_s(\ckG)}$ 在求逆下不变。这蕴含 $\sum_i\ell_i=r_s(\ckG)h/2$。故 $\dim V_{\theta^\vee}=\sum_i(2\ell_i+1)=(h+1)r_s(\ckG)$。由于 $V_{\theta^\vee}$ 的非零权空间以 $\ckG$ 的短根为指标，从而以轨道 $W\theta^\vee$ 为指标，得 $\#(W\theta^\vee)=\dim V_{\theta^\vee}-\dim V_{\theta^\vee}(0)=hr_s(\ckG)$。这证明 (4) 与其余各数一致。
\end{proof}

\section{附录：$G_2$ 的伴随 Schubert 簇}

本节假设 $\operatorname{char}(k)>3$。回忆（见 [8] 中的表）$G_2$ 有四个非正则幂幺轨道：
\begin{itemize}[leftmargin=1.6em, itemsep=.2em]
\item 次正则轨道，含 $P_{-\theta}$ 的幂幺根中的一般元；其维数为 $10$；
\item 含 $U^\times_{-\gamma}$ 的轨道，维数为 $8$；
\item 含 $U^\times_{-\theta}$ 的轨道，维数为 $6$；
\item 恒等轨道，维数为 $0$。
\end{itemize}
对 $G_2$，$\mathrm{Gr}^{\mathrm{triv}}_{\le\gamma^\vee}$ 的 $G$-轨道与这些幂幺轨道密切相关。更精确地：

\begin{lemma}
\begin{enumerate}[label=(\arabic*), leftmargin=1.9em, itemsep=.25em]
\item $\mathrm{Gr}^{\mathrm{triv}}_{\le\gamma^\vee}$ 上有四个 $\operatorname{Ad}(G)$-轨道
\[
\mathrm{Gr}^{\mathrm{triv}}_{\le\gamma^\vee}=\mathrm{Gr}^{\mathrm{subr}}_G\sqcup\operatorname{Ad}(G)U^\times_{-\gamma,-1}\sqcup\mathrm{Gr}^{\mathrm{triv}}_{\theta^\vee}\sqcup\{\star\},
\]
其维数分别为 $10,8,6,0$；在赋值映射 $\mathrm{ev}_{\tau=1}$ 下它们映到上述四个非正则幂幺轨道。
\item 态射
\[
\begin{aligned}
\nu:\ \frac{G}{P_{-\theta}}\times\Bigl(\prod_{\langle\beta,-\theta^\vee\rangle\ge1}U_{\beta,-1}\Bigr)&\longrightarrow\mathrm{Gr}^{\mathrm{triv}}_{\le\gamma^\vee},\\
\Bigl(g,\prod u_\beta(c_\beta\tau^{-1})\Bigr)&\mapsto\operatorname{Ad}(g)\Bigl(\prod u_\beta(c_\beta\tau^{-1})\Bigr)
\end{aligned}
\]
是消解。它在各 $G$-轨道上的纤维为：
\begin{itemize}[leftmargin=1.4em, itemsep=.1em]
\item $\nu$ 在 $\mathrm{Gr}^{\mathrm{subr}}$ 上是同构；
\item $\nu^{-1}(u_{-\gamma}(\tau^{-1}))\cong\PP^1$；
\item $\nu^{-1}(u_{-\theta}(\tau^{-1}))$ 是 $\PP^1$ 上的射影锥（它含一个点，其补是 $\PP^1$ 上的线丛）；
\item $\nu^{-1}(\star)\cong G/P_{-\theta}$。
\end{itemize}
\item 按与 $\nu$ 相同方式定义的态射
\[
\nu':\ \frac{G}{P_{-\gamma}}\times\Bigl(\prod_{\langle\beta,-\gamma^\vee\rangle\ge2}U_{\beta,-1}\Bigr)\longrightarrow\mathrm{Gr}^{\mathrm{triv}}_{\le\gamma^\vee}
\]
是 $\operatorname{Ad}(G)U^\times_{-\gamma,-1}$ 的闭包的消解。它在各 $G$-轨道上的纤维为：
\begin{itemize}[leftmargin=1.4em, itemsep=.1em]
\item $\nu'$ 在 $\operatorname{Ad}(G)U^\times_{-\gamma,-1}$ 上是同构；
\item $\nu'^{-1}(u_{-\theta}(\tau^{-1}))\cong\PP^1$；
\item $\nu'^{-1}(\star)\cong G/P_{-\gamma}$。
\end{itemize}
\end{enumerate}
\end{lemma}

\noindent\textbf{对 $G=G_2$ 证明定理 4(2)。} 经与其他型情形相同的化归步骤，只需证
\begin{equation}
\chi_c\bigl(\mathrm{Gr}^{\mathrm{triv},a=0}_{\gamma^\vee},\mathrm{ev}^*J\bigr)=-r_s(G)=-1.
\end{equation}

\medskip\noindent\textbf{第一步.} $\chi_c\bigl((\operatorname{Ad}(G)U^\times_{-\gamma,-1})^{a=0},\mathrm{ev}^*J\bigr)=-1$。

对任何短根 $\beta$，令 $V_\beta=\prod_{\langle\beta',\beta^\vee\rangle\ge2}U_{\beta',-1}$。则 $\nu'$ 的源有 Bruhat 分解
\[
\frac{G}{P_{-\gamma}}\times V_{-\gamma}=\bigsqcup_{\beta\ \text{短根}}\operatorname{Ad}(U)V_\beta.
\]
下述断言可仿 §5.6 中的断言证明。

\medskip\noindent\textbf{断言.} 对短根 $\beta$，
\[
\chi_c\bigl(\operatorname{Ad}(U)V^{>-\theta}_\beta,\mathrm{ev}^*J\bigr)=
\begin{cases}
0, & \Phi^\beta_{\ge2}\ \text{含单根},\\
1, & \text{否则}.
\end{cases}
\]

观察根系 $G_2$，有 $3$ 个短根 $\beta$ 使 $\Phi^\beta_{\ge2}$ 不含单根。因此 $\chi_c\bigl((\frac{G}{P_{-\gamma}}\times V_{-\gamma})^{a=0},\nu'^*\mathrm{ev}^*J\bigr)=3$。另一方面，由引理 11.1(3) 有
\[
\begin{aligned}
3&=\chi_c\Bigl(\bigl(\tfrac{G}{P_{-\gamma}}\times V_{-\gamma}\bigr)^{a=0},\nu'^*\mathrm{ev}^*J\Bigr)\\
&=\chi_c(G/P_{-\gamma})+\chi_c(\PP^1)\chi_c\bigl(\mathrm{Gr}^{\mathrm{triv},a=0}_{\theta^\vee},\mathrm{ev}^*J\bigr)\\
&\quad+\chi_c\bigl((\operatorname{Ad}(G)U^\times_{-\gamma,-1})^{a=0},\mathrm{ev}^*J\bigr).
\end{aligned}
\]
代入 $\chi_c(G/P_{-\gamma})=\#W/W_{\gamma^\vee}=6$、$\chi_c(\PP^1)=2$ 以及由定理 4(1) 的证明得到的 $\chi_c(\mathrm{Gr}^{\mathrm{triv},a=0}_{\theta^\vee},\mathrm{ev}^*J)=-1$，即得 $\chi_c\bigl((\operatorname{Ad}(G)U^\times_{-\gamma,-1})^{a=0},\mathrm{ev}^*J\bigr)=-1$。

\medskip\noindent\textbf{第二步.} $\chi_c\bigl(\mathrm{Gr}^{a=0}_{\mathrm{subr}},\mathrm{ev}^*J\bigr)=0$。

对任何长根 $\alpha$，令 $V_\alpha=\prod_{\langle\alpha',\alpha^\vee\rangle\ge1}U_{\alpha',-1}$。则 $\nu$ 的源有 Bruhat 分解
\[
\frac{G}{P_{-\theta}}\times V_{-\theta}=\bigsqcup_{\alpha\ \text{长根}}\operatorname{Ad}(U)V_\alpha.
\]
下述断言可仿 §5.6 中的断言证明。

\medskip\noindent\textbf{断言.} 对长根 $\alpha$，
\[
\chi_c\bigl(\operatorname{Ad}(U)V^{>-\theta}_\alpha,\mathrm{ev}^*J\bigr)=
\begin{cases}
0, & \Phi^\alpha_{\ge1}\ \text{含单根},\\
1, & \text{否则}.
\end{cases}
\]

观察根系 $G_2$，$\alpha=-\theta$ 是唯一使 $\Phi^\alpha_{\ge1}$ 不含单根的长根。因此 $\chi_c\bigl((\frac{G}{P_{-\theta}}\times V_{-\theta})^{a=0},\nu^*\mathrm{ev}^*J\bigr)=1$。另一方面，由引理 11.1(2) 有
\begin{align*}
1&=\chi_c\Bigl(\bigl(\tfrac{G}{P_{-\theta}}\times V_{-\theta}\bigr)^{a=0},\nu^*\mathrm{ev}^*J\Bigr)\\
&=\chi_c(G/P_{-\theta})+\chi_c(\PP^1)\chi_c\bigl((\operatorname{Ad}(G)U^\times_{-\gamma,-1})^{a=0},\mathrm{ev}^*J\bigr)\\
&\qquad+\chi_c\bigl(\nu^{-1}(u_{-\theta}(\tau^{-1}))\bigr)\chi_c\bigl(\mathrm{Gr}^{\mathrm{triv},a=0}_{\theta^\vee},\mathrm{ev}^*J\bigr)
+\chi_c\bigl(\mathrm{Gr}^{a=0}_{\mathrm{subr}},\mathrm{ev}^*J\bigr).
\end{align*}
代入 $\chi_c(G/P_{-\theta})=\#W/W_\theta=6$、$\chi_c(\PP^1)=2$、由引理 11.1(2) 得到的 $\chi_c(\nu^{-1}(u_{-\theta}(\tau^{-1})))=3$、第一步的 $\chi_c\bigl((\operatorname{Ad}(G)U^\times_{-\gamma,-1})^{a=0},\mathrm{ev}^*J\bigr)=-1$ 以及定理 4(1) 的证明中的 $\chi_c(\mathrm{Gr}^{\mathrm{triv},a=0}_{\theta^\vee},\mathrm{ev}^*J)=-1$，得 $\chi_c\bigl(\mathrm{Gr}^{a=0}_{\mathrm{subr}},\mathrm{ev}^*J\bigr)=0$。\zhnote{原文此处写作 ``we conclude that $\chi_c(\mathrm{Gr}^{a=0}_{\mathrm{subr}},\mathrm{ev}^*J)=-1$''，与其小标题（$=0$）及上式代入的结果（$6-2-3+x=1$ 给出 $x=0$）均不符；且只有取 $0$ 才能在末尾与第一步的 $-1$ 相加得到 (11.2)。故按 $0$ 译。[来源：原文 p.55 第 2 步末尾与 (11.2)]}

合并第一步与第二步，由 $\mathrm{Gr}^{\mathrm{triv},a=0}_{\gamma^\vee}=\mathrm{Gr}^{a=0}_{\mathrm{subr}}\sqcup(\operatorname{Ad}(G)U^\times_{-\gamma,-1})^{a=0}$ 即得 (11.2)。

这证明了 $G_2$ 情形的定理 4(2)。
